\section[Multipliers and Global Positivity]{Multipliers, Jacobi Identities, and Global Positivity}
\label{sec:algebraic-jacobi}

The multiplier expansion relates global divergence to its Jacobi trace;
\cref{sec:laplacian-calculus} reduces the sign test to the invariant image.

\subsection{The global multiplier defect}

\begin{proposition}[Odd-power invariance and regular defining potentials]
\label{prop:potential-normal-form}
Let \(\Omega\subset\R^N\setminus\{0\}\) be open and conical, and
let \(F\) be smooth and homogeneous with \(\nabla F\neq0\) on
\(C\cap\Omega=\{F=0\}\cap\Omega\).
\begin{enumerate}[label=\textup{(\roman*)}]
 \item Let \(a>0\) be smooth and homogeneous on \(\Omega\), and
 let \(m\geq1\) be an odd integer.  Set
 \[
  \widetilde F=F^m,\qquad G=aF,\qquad
  \widetilde G=a^m\widetilde F=G^m.
 \]
 Then
 \begin{equation}\label{eq:odd-power-field}
  \nabla\widetilde G=mG^{m-1}\nabla G,\qquad
  \cL(\widetilde G)=m^3G^{3m-3}\cL(G).
 \end{equation}
 Thus \(X_{\widetilde G}=X_G\) wherever \(F\neq0\) and
 \(\nabla G\neq0\), and \(X_G\) gives its smooth extension across
 \(C\cap\Omega\).  Conversely, every positive smooth homogeneous
 multiplier \(b\) for \(\widetilde F\) is obtained by setting
 \(a=b^{1/m}\).  Hence the two defining equations admit the same
 multiplier subcalibration fields, with the same regular boundary
 extensions and the same remaining exceptional sets.
 \item If a smooth homogeneous function \(G\) has the same zero set
 and signs as \(F\) on \(\Omega\), and \(\nabla G\neq0\) on
 \(C\cap\Omega\), then \(G=bF\) for a positive smooth homogeneous
 function \(b\) on \(\Omega\).
\end{enumerate}
\end{proposition}

\begin{proof}
For \textup{(i)}, the chain rule gives \(\cL(\psi(G))=(\psi'(G))^3\cL(G)\).
Take \(\psi(t)=t^m\).
Its derivative is positive off \(t=0\), which proves the field
identity.  On the regular cone, \(\nabla G=a\nabla F\neq0\),
so \(X_G\) provides the exterior normal extension.  Off the cone,
the two gradients vanish at exactly the same points.  The inverse
correspondence follows from \(bF^m=(b^{1/m}F)^m\); the positive
root is smooth and homogeneous.

For \textup{(ii)}, use local coordinates \((y,t)\) with \(t=F\).
Since \(G(y,0)=0\),
\[
 G(y,t)=t\int_0^1\partial_tG(y,st)\,ds.
\]
The integral extends \(b=G/F\) smoothly, with
\(b|_C=|\nabla G|/|\nabla F|>0\) by the common phase signs and
regularity.  These extensions agree on overlaps; the quotient also
gives \(\deg b=\deg G-\deg F\).
\end{proof}

Part \textup{(i)} preserves the divergence and its phase integrals.
For \(m>1\), \(\nabla\widetilde G=0\) on the cone; the
Jacobi--Hardy formulas use its regular field extension \(X_G\), the degree
\(\ell=\deg\widetilde G/m\), and \(w=|\nabla G|^{-1}|_C\).

The displayed algebra \(\R[R,F^m]\) need not itself be Laplacian closed, even when
\(\R[R,F]\) is: for \(F=x_1\),
\(\Delta(F^3)=6F\notin\R[R,F^3]\).
The radial and gradient powers used below are discussed in
\cref{app:power-identities}.

Let \(\Omega\subset\R^N\) be open and \(F,a\in C^\infty(\Omega)\),
with \(a>0\). Suppose \(\nabla F\ne0\) on
\(C=\{F=0\}\cap\Omega\) and \(C\) is minimal. Put \(G=aF\).
The product rule gives \(\cL(G)|_C=a^3\cL(F)|_C=0\).
Smooth division by \(F\), as in
\cref{prop:potential-normal-form}\textup{(ii)}, therefore defines
\[
 \lambda_F=\frac{\cL(F)}F,\qquad \cK_F(a)=\frac{\cL(aF)}F.
\]
Then \(G\cL(G)=aF^2\cK_F(a)\). Polynomial divisibility gives
the same \(\lambda_F\) as \eqref{eq:L-divisibility}. Write
\[
 P=\nabla F,\quad Q=\nabla a,\qquad
 S=|P|^2,\quad T=\langle P,Q\rangle,\quad V=|Q|^2.
\]

\begin{theorem}[Exact multiplier expansion]
\label{thm:exact-multiplier}
Under the preceding smooth hypotheses, the multiplier defect equals
\begin{align}
\cK_F(a)={}&a^3\lambda_F\notag\\
&+a^2\bigl(S\Delta a+2T\Delta F-2\Hess F(P,Q)-\Hess a(P,P)\bigr)\notag\\
&+2a(T^2-SV)\notag\\
&+Fa\bigl(2T\Delta a+V\Delta F-\Hess F(Q,Q)
              -2\Hess a(P,Q)\bigr)\notag\\
&+F^2\cL(a).                                           \label{eq:exact-multiplier}
\end{align}
If, in addition, \(F,a\) are polynomials in a Laplacian algebra
\(\cA\) and \(F\mid\cL(F)\) in \(\R[x_1,\ldots,x_N]\),
then every term in \eqref{eq:exact-multiplier} belongs to \(\cA\).
\Cref{prop:divisibility} gives sufficient geometric conditions for
this polynomial divisibility.
\end{theorem}

\begin{proof}
Expand \(\nabla(aF)=aP+FQ\),
\(\Delta(aF)=a\Delta F+F\Delta a+2T\), and
\[
 \Hess(aF)=a\Hess F+F\Hess a+dF\otimes da+da\otimes dF.
\]
The constant term in \(\cL(aF)\) is
\(a^3\cL(F)=a^3F\lambda_F\); collecting the remaining powers of \(F\)
gives the formula.  The identity
\[
 2\Hess u(\nabla v,\nabla w)
 =\Gamma(v,\Gamma(u,w))+\Gamma(w,\Gamma(u,v))
  -\Gamma(u,\Gamma(v,w))
\]
expresses the Hessian terms through gradient inner products.

By closure, \(\cL(F)=S\Delta F-\Gamma(F,S)/2\in\cA\).
Polynomial division, \cref{prop:MR-finite-and-division}, then gives
\(\lambda_F\in\cA\), proving the algebra assertion.
\end{proof}

\begin{lemma}[The defect at a critical zero]
\label{lem:critical-multiplier-defect}
Let \(F\) be a nonzero polynomial satisfying
\(\cL(F)=F\lambda_F\), with polynomial \(\lambda_F\), and let
\(a\) be a \(C^2\) function near a point \(p\) with
\(F(p)=0\) and \(\nabla F(p)=0\).
The right-hand side of \eqref{eq:exact-multiplier} defines a continuous
extension of \(\cL(aF)/F\), and
\begin{equation}\label{eq:critical-multiplier-defect}
 \cK_F(a)(p)=a(p)^3\lambda_F(p).
\end{equation}
If \(\lambda_F(p)<0\), no such multiplier with \(a(p)>0\)
can satisfy \((aF)\cL(aF)\ge0\) throughout a punctured neighborhood
of \(p\). In particular, a nonnegative \(C^2\) multiplier satisfying
this sign condition must vanish at \(p\).
\end{lemma}

\begin{proof}
The expansion is a product-rule identity wherever \(F\ne0\), and
every term on its right-hand side is continuous. At \(p\), one has
\(F=0\), \(P=0\), \(S=0\), and \(T=0\), so only
\(a^3\lambda_F\) remains. If \(a(p)>0\) and
\(\lambda_F(p)<0\), continuity gives \(a>0\) and
\(\cK_F(a)<0\) near \(p\). A nonzero polynomial has a zero set
with empty interior, so there are nearby points with \(F\ne0\), where
\[
 (aF)\cL(aF)=aF^2\cK_F(a)<0.
\]
This contradicts the required sign.
\end{proof}

\paragraph{Logarithmic differentiation of a positive multiplier.}
On the open domain where \(a>0\), including the regular cone,
\(h=\log a\) is smooth and satisfies
\begin{equation}\label{eq:log-multiplier-derivatives}
 \nabla h=\frac{\nabla a}{a},\qquad
 \Delta h=\frac{\Delta a}{a}-\frac{|\nabla a|^2}{a^2},\qquad
 \Gamma(h,q)=\frac{\Gamma(a,q)}a.
\end{equation}
If $a$ has degree $d$, then $h(tx)=h(x)+d\log t$ for $t>0$.

\subsection{The Jacobi boundary condition}
\label{sec:jacobi-boundary-condition}

Let \(J_{C_{\rm reg}}=\Delta_{C_{\rm reg}}+|A_C|^2\)
be the Jacobi operator associated with the second variation; see
\cite{Simons68} and \cite[Remark~9.8]{Simon83}. With scalar mean curvature
defined by \(H=\diver_{C_{\rm reg}}\nu\), a variation with normal
speed \(w\) has linearization \(-J_Cw\).

\begin{theorem}[Exact Jacobi trace]
\label{thm:Jacobi-trace}
For the smooth \(F,a\) of \cref{thm:exact-multiplier}, along
\(C_{\reg}\cap\Omega\),
\begin{equation}\label{eq:Jacobi-trace}
 \boxed{
 \left.\cK_F(a)\right|_{C_{\rm reg}}
 =-a^4|\nabla F|^3J_{C_{\rm reg}}
 \left(\frac1{a|\nabla F|}\right).}
\end{equation}
Consequently, \(\cK_F(a)\geq0\) near this hypersurface requires the positive
Jacobi supersolution
\begin{equation}\label{eq:positive-Jacobi}
 w=(a|\nabla F|)^{-1}>0,
 \qquad J_Cw\leq0.
\end{equation}
It is a Jacobi field exactly when the boundary defect vanishes.
\end{theorem}

\begin{proof}
The flow of \(\nabla G/|\nabla G|^2\) parametrizes the level sets
\(M_t=\{G=t\}\), with normal speed
\(w=|\nabla G|^{-1}=(a|\nabla F|)^{-1}\) at \(t=0\).
The mean-curvature variation at the minimal zero level is
\(\frac{d}{dt}H_{M_t}|_0=-J_Cw\).
On the other hand, \eqref{eq:div-normalized-gradient} gives, along this flow,
\[
 H_{M_t}=\frac{F\cK_F(a)}{|\nabla G|^3},\qquad
 \frac{dF}{dt}\bigg|_0
 =\left.\frac{\langle\nabla F,\nabla G\rangle}{|\nabla G|^2}\right|_C
 =\frac1a.
\]
Since \(F=0\) on \(C\), differentiation of the other factors contributes
zero at \(t=0\).  Using \(|\nabla G|^3=a^3|\nabla F|^3\) there yields
\[
 -J_Cw=\frac{d}{dt}H_{M_t}\bigg|_0
 =\left.\frac{\cK_F(a)}{a^4|\nabla F|^3}\right|_C,
\]
which is \eqref{eq:Jacobi-trace}.
\end{proof}

\begin{proposition}[Harmonic power multipliers on the boundary]
\label{prop:harmonic-power-Jacobi}
Let \(F\not\equiv0\) be a homogeneous polynomial of degree \(e\ge2\) satisfying
\eqref{eq:harmonic-minimal-weight}. Put \(R=r^2\),
\(K=\Delta S=2|\Hess F|^2\), and \(T=|\nabla S|^2\).
For real \(\sigma,\gamma\), define \(G=r^\sigma S^\gamma F\) on
\(\{S>0\}\). Its boundary speed is \(w=r^{-\sigma}S^{-\gamma-1/2}\).
On \(C_{\reg}\),
\begin{equation}\label{eq:harmonic-power-Jacobi}
 -\frac{J_Cw}{w}
 =\frac{\sigma\bigl(N-1-2e-\sigma-4\gamma(e-1)\bigr)}{R}
  +\frac{\gamma K-(2\gamma+1)q}{S}
  -\frac{\gamma(2\gamma+3)}2\frac{T}{S^2}.
\end{equation}
In particular, a two-sided subcalibration generated by \(G\) requires
the right-hand side to be nonnegative on \(C_{\reg}\).
\end{proposition}

\begin{proof}
Set \(w_0=S^{-1/2}\). Applying \cref{thm:Jacobi-trace} with
\(a=1\) and \(\cL(F)=-qF\) gives \(J_Cw_0/w_0=q/S\).
On \(C_{\reg}\), equation \eqref{eq:harmonic-minimal-weight} implies
\(\nabla S\perp\nabla F\) and \(\Hess F(\nu_C,\nu_C)=0\).
Since \(\Hess F\,\nu_C=\nabla S/(2\sqrt S)\), decomposition into
normal and tangential components gives
\[
 |A_C|^2=\frac{K}{2S}-\frac{T}{2S^2},\qquad
 |\nabla_C\log w_0|^2=\frac{T}{4S^2}.
\]
For \(p=2\gamma+1\), the chain rule reads
\[
 \frac{J_C(w_0^p)}{w_0^p}
 =p\frac{J_Cw_0}{w_0}+(1-p)|A_C|^2
       +p(p-1)|\nabla_C\log w_0|^2.
\]
Finally, \(C\) has dimension \(N-1\), so
\(\Delta_Cr^{-\sigma}/r^{-\sigma}
 =\sigma(\sigma-N+3)/R\).
The homogeneity of \(w_0^p\), of degree \(-p(e-1)\), gives the
cross term \(2\sigma p(e-1)/R\). Substitution proves the formula.
\end{proof}

\begin{theorem}[Positive-solution identity and homogeneous Hardy gap]
\label{thm:ground-state-gap}
Let \(C^n\subset\R^{n+1}\), \(n\geq2\), be a cone whose regular part
is two-sided and minimal.  Let \(G\) be a smooth homogeneous defining function of degree
\(\ell\) on a conical neighborhood of \(C_{\reg}\setminus\{0\}\), with
\(\nabla G\neq0\) there.  Set
\begin{equation}\label{eq:boundary-Jacobi-potential}
 w=|\nabla G|^{-1}\big|_C>0,\qquad
 V_G=-\frac{J_Cw}{w}=r^{-2}\beta_G(\theta).
\end{equation}
For every \(u\in C_c^\infty(C_{\reg}\setminus\{0\})\),
\begin{equation}\label{eq:ground-state-identity}
 Q_C(u)=\int_Cw^2\left|\nabla_C\left(\frac uw\right)\right|^2
             +\int_CV_Gu^2.
\end{equation}
If \(\beta_G\geq\beta_0\) on the regular link for a finite constant \(\beta_0\), then
\begin{equation}\label{eq:homogeneous-strict-stability}
 Q_C(u)\geq\left[\left(\ell-\frac n2\right)^2+\beta_0\right]
                  \int_Cr^{-2}u^2.
\end{equation}
For \(G=aF\) in \cref{thm:Jacobi-trace},
\begin{equation}\label{eq:boundary-potential-algebraic}
 V_G=\left.\frac{\cK_F(a)}{a^3|\nabla F|^2}\right|_C
     =\left.\partial_{\nu_C}\diver
          \left(\frac{\nabla G}{|\nabla G|}\right)\right|_C.
\end{equation}
Consequently a nonnegative Jacobi boundary defect implies stability;
if also \(\ell\neq n/2\), it implies strict stability with constant at
least \((\ell-n/2)^2\).
\end{theorem}

\begin{proof}
Substituting \(u=wv\) and integrating by parts gives
\eqref{eq:ground-state-identity}.
Homogeneity gives \(w=r^{1-\ell}\phi(\theta)\), \(\phi>0\).
Writing \(\Lambda_{\reg}=C_{\reg}\cap\Sph^n\), the cone volume
element \(r^{n-1}\,dr\,d\mu_\Lambda\) gives
\[
 \int_Cw^2|\nabla_Cv|^2
 \geq\int_{\Lambda_{\reg}}\phi^2
       \int_0^\infty r^{n+1-2\ell}|\partial_rv|^2\,dr\,d\mu_\Lambda.
\]
For every real \(p\) and \(f\in C_c^\infty(0,\infty)\), completion of the
square proves
\[
 \int_0^\infty r^p(f')^2\,dr
 \geq\frac{(p-1)^2}{4}\int_0^\infty r^{p-2}f^2\,dr.
\]
Apply this with \(p=n+1-2\ell\) for each angular point.
The resulting term is \((\ell-n/2)^2\int_Cr^{-2}u^2\).
Adding \(\int_CV_Gu^2\geq\beta_0\int_Cr^{-2}u^2\) proves
\eqref{eq:homogeneous-strict-stability}.
The first equality in \eqref{eq:boundary-potential-algebraic} follows by
dividing \eqref{eq:Jacobi-trace} by \(a^4|\nabla F|^3w\).
For the second, differentiate
\(\diver X_G=F\cK_F(a)/|\nabla G|^3\) in the unit normal direction on
\(F=0\), where \(\partial_\nu F=|\nabla F|\).
\end{proof}

\begin{remark}[Vanishing trace and the Hardy gap]
\label{rem:vanishing-trace-strict-stability}
If \(J_Cw=0\), formula \eqref{eq:ground-state-identity} still contains
the Hardy gap \((\ell-n/2)^2\).  This lower bound vanishes only at the central degree
\(\ell=n/2\).
Failure of strict stability requires a sequence with Hardy quotient
tending to zero; see \cref{prop:complex-zero-gap}.
\end{remark}

\begin{remark}[Subcalibrations as paired geometric barriers]
\label{rem:paired-barriers}
For \(G=aF\), \(a>0\), and \(\cK_F(a)\geq0\), the regular levels
\(M_t=\{G=t\}\), oriented by \(X_G\), satisfy
\[
 H_{M_t}=\frac{F\cK_F(a)}{|\nabla G|^3},\qquad
 tH_{M_t}\geq0.
\]
Thus \(G\) is a supersolution on \(F<0\) and a subsolution on
\(F>0\) of the level-set equation \(\diver X_G=0\).
The level-set variation in the proof of \cref{thm:Jacobi-trace}
has normal graph height \(s_t=tw+O(t^2)\), with
\(w=(a|\nabla F|)^{-1}\), and gives
\begin{equation}\label{eq:paired-barrier-expansion}
 H_{M_t}=-tJ_Cw+O(t^2).
\end{equation}
Thus \(J_Cw<0\) gives strict local barriers on both sides.
\end{remark}

The sign of \(\cK_F(a)\) away from the cone requires a separate
calculation, as the stable non-minimizing examples in
\cref{sec:isoparametric-reduction} illustrate.

\subsubsection*{The admissible degree}

Let \(C^{N-1}\subset\R^N\), \(N\geq3\), be a minimal cone, let
\(\Sigma=C\cap\Sph^{N-1}\), and suppose a homogeneous potential \(G=aF\)
has total degree \(\ell\).  Then \(w=(a|\nabla F|)^{-1}\) has degree
\(1-\ell\).  If the entire link \(\Sigma\) is smooth, it is compact
and connected; see the maximum-principle argument of Hardt--Simon
\cite[p.~104]{HardtSimon85}.
Let \(\lambda_1(\Sigma)\) be the first eigenvalue of
\(-\Delta_\Sigma-|A_\Sigma|^2\).  Existence of
a positive homogeneous Jacobi supersolution requires
\begin{equation}\label{eq:degree-window}
 \ell\in I_F:=
 \left[\frac{N-1}{2}-\beta_F,
       \frac{N-1}{2}+\beta_F\right],
 \qquad
 \beta_F^2=\frac{(N-3)^2}{4}+\lambda_1(\Sigma).
\end{equation}
If the radicand is negative, no positive homogeneous Jacobi supersolution
with positive angular part exists.  More precisely, if \(\phi_1>0\) is the
first eigenfunction and
\(w=r^{1-\ell}\phi_1\), then
\[
 J_Cw=r^{-1-\ell}\left[
 (1-\ell)(N-2-\ell)-\lambda_1(\Sigma)
 \right]\phi_1.
\]
For any positive angular supersolution, integrating its inequality
against \(\phi_1\) gives the same necessary degree condition.
Using \(\phi_1\), interior degrees give \(J_Cw<0\) and endpoints
give \(J_Cw=0\). A prescribed angular multiplier requires its own sign test.

For a smooth compact connected link, the zero-gap case is rigid.
Write $n=N-1$, $H_\Sigma=-\Delta_\Sigma-|A_\Sigma|^2$, and
$\delta_C=(n-2)^2/4+\lambda_1(\Sigma)$, the exact Hardy gap in
\eqref{eq:exact-Hardy-gap}. If $w=r^{1-\ell}\phi>0$ and
$-J_Cw/w=r^{-2}\beta\ge0$, then
\begin{equation}\label{eq:critical-boundary-rigidity}
 \delta_C=0\quad\Longrightarrow\quad
 \ell=n/2,\qquad \beta\equiv0,\qquad
 \phi\text{ is a positive first eigenfunction of }H_\Sigma.
\end{equation}
Indeed, \eqref{eq:degree-window} forces $\ell=n/2$.
Then $(H_\Sigma+(n-2)^2/4)\phi=\beta\phi$.
Pairing with a positive first eigenfunction $\psi$ gives
$0=\int_\Sigma\beta\phi\psi$, hence $\beta=0$.
Simplicity of the first eigenvalue now gives the last assertion.
In particular, a nonnegative boundary defect that is positive at one
point suffices for strict stability on such a link.

\begin{remark}[Two admissible power multipliers]
\label{cor:two-power-strict-stability}
Let $F$ be a homogeneous polynomial of degree $e\ge2$ on
$\R^{n+1}$, $n\ge3$, whose zero set $C$ is minimal, with
$\nabla F\ne0$ on $C\setminus\{0\}$ and nonempty smooth compact
connected link $\Sigma$. Put $S=|\nabla F|^2$, and suppose that
$s=S|_\Sigma$ is nonconstant. If two distinct real parameter pairs
$(\alpha_i,b_i)$ give potentials
\[
 G_i=r^{\alpha_i}S^{b_i/2}F,\qquad
 \operatorname{sgn}(F)\diver X_{G_i}\ge0\quad\text{where }F\ne0
\]
in a two-sided neighborhood of $C\setminus\{0\}$ on which both
fields are defined, then $C$ is strictly stable.
Indeed, by \cref{thm:Jacobi-trace}, the positive boundary speeds
$w_i=r^{-\alpha_i}S^{-(b_i+1)/2}$ are Jacobi supersolutions.
\Cref{thm:ground-state-gap} first gives stability.
Their potential degrees are $\ell_i=e+\alpha_i+b_i(e-1)$.
If the Hardy gap were zero, \eqref{eq:critical-boundary-rigidity}
would force $\ell_i=n/2$ and both angular parts
$s^{-(b_i+1)/2}$ to be proportional to the same first eigenfunction.
Thus $s^{-(b_1-b_2)/2}$ would be constant. Since $s$ is nonconstant,
$b_1=b_2$, and equality of the degrees then gives $\alpha_1=\alpha_2$,
a contradiction.
\end{remark}

\paragraph{Example: admissible degrees on the Simons cone.}
\label{par:Simons-degree-example}
Let
\[
 F(x,y)=|x|^2-|y|^2,\qquad (x,y)\in\R^4\times\R^4,\qquad R=r^2.
\]
Its zero link is
\(\Sigma=\Sph^3(1/\sqrt2)\times\Sph^3(1/\sqrt2)\), with
\(|A_\Sigma|^2=6\).  The constant function attains
\(\lambda_1(\Sigma)=-6\), so \eqref{eq:degree-window} gives
\[
 \beta_F^2=\frac{25}{4}-6=\frac14,\qquad I_F=[3,4].
\]
The quadratic \(F\) itself lies outside this interval.  Choose instead
\(G=RF=|x|^4-|y|^4\), of degree \(\ell=4\), as in
\cite[Section~2]{DePhilippisPaolini}; see
\cref{rem:subcalibration-history} for its historical antecedents.
Direct calculation gives
\begin{align*}
 |\nabla G|^2&=4R(R^2+3F^2),& \Delta G&=24F,\\
 \cL(G)&=192RF^3,& \cK_F(R)&=192RF^2,\\
 G\cL(G)&=192R^2F^4\geq0.
\end{align*}
This is the strictly minimizing \(g=2\), \(m_1=m_2=3\) case of
\cref{sec:isoparametric-reduction}.

On the other hand,
\[
 w=|\nabla G|^{-1}\big|_C=\frac1{2r^3},\qquad
 J_Cw=0,\qquad \cK_F(R)\big|_C=0.
\]
Nevertheless \cref{thm:ground-state-gap} gives
\[
 Q_C(u)\geq\left(4-\frac72\right)^2\int_Cr^{-2}u^2
 =\frac14\int_Cr^{-2}u^2,
\]
and \eqref{eq:exact-Hardy-gap} shows this constant is optimal.
Thus an endpoint degree and zero Jacobi trace coexist with both
strict area minimality and strict stability.
For \(G_\eps=r^{-\eps}G\), \(0<\eps<1\), the boundary coefficient
is \(\beta_{G_\eps}=\eps-\eps^2>0\).
For all sufficiently small \(\eps>0\), \(G_\eps\) retains the
subcalibration sign and its normalized gradient satisfies the linear
distance bound; see \cref{cor:four-degenerate-certificates}.

\subsubsection*{The spectral obstruction for a singular link}
For a singular link, we use the angular Rayleigh infimum, without
assuming connectedness of \(\Sigma_{\reg}\) or existence of a first
eigenfunction.\footnote{If \(\Sigma\) is the connected support of a stationary
integral \((N-2)\)-varifold in \(\Sph^{N-1}\), with
\(\cH^{N-3}(\Sing\Sigma)=0\), regular-part connectedness follows from
\cite[Lemma~2.2]{Zhu18}.}
Define
\begin{equation}\label{eq:Friedrichs-bottom}
 \lambda_{\rm F}(\Sigma)=
 \inf_{0\neq\phi\in C_c^\infty(\Sigma_{\rm reg})}
 \frac{\int_{\Sigma_{\rm reg}}
 (|\nabla\phi|^2-|A_\Sigma|^2\phi^2)}
 {\int_{\Sigma_{\rm reg}}\phi^2}.
\end{equation}
This infimum may equal \(-\infty\); when finite, it is the bottom of
the Friedrichs extension and need not be attained.
Zhu's spherical stability convention
\cite[Section~2.1.5, equation~(2.8)]{Zhu18} includes the Ricci term:
\[
 \lambda_1^{Z}(\Sigma)=\lambda_{\rm F}(\Sigma)-(N-2).
\]
The exact cone stability quotient is
\begin{equation}\label{eq:exact-Hardy-gap}
\begin{aligned}
 \inf_{0\neq u\in C_c^\infty(C_{\rm reg}\setminus\{0\})}
 \frac{Q_C(u)}{\int_Cr^{-2}u^2}
 &=\frac{(N-3)^2}{4}+\lambda_{\rm F}(\Sigma)\\
 &=\frac{(N-1)^2}{4}+\lambda_1^{Z}(\Sigma).
\end{aligned}
\end{equation}
By radial--angular separation \cite{Simons68,CabrePoggesi18},
the sharp radial Hardy inequality and the angular Rayleigh bound give
the lower bound. For a product
test \(u(r,\theta)=\eta(r)\varphi(\theta)\), the cone quotient is
the sum of the radial Hardy quotient and the angular Rayleigh
quotient. Independent minimizing sequences give equality, including
\(\lambda_{\rm F}(\Sigma)=-\infty\).
If the right-hand side is negative, no fields satisfying the
comparison criterion can exist.

\begin{remark}[What the boundary trace determines]
Equality \(a_1|_{C_{\rm reg}}=a_2|_{C_{\rm reg}}\) is equivalent to
local smooth divisibility of \(a_1-a_2\) by \(F\). By
\eqref{eq:Jacobi-trace}, it implies
\begin{equation}\label{eq:boundary-ideal-class}
 \left.\cK_F(a_1)\right|_{C_{\rm reg}}
 =\left.\cK_F(a_2)\right|_{C_{\rm reg}}.
\end{equation}
Thus the boundary defect depends only on \(a|_C\).
\end{remark}

For calibrated complex cones with zero Hardy gap, see the
Jacobi-field construction in \cref{prop:complex-zero-gap}.

\subsection{A multiplier criterion for harmonic defining functions}
\label{sec:harmonic-multiplier-criterion}

The following criterion permits multipliers that degenerate at the
singular set. Its spectral application appears in
\cref{sec:spectral-positivity}.

Let \(F\not\equiv0\) be a homogeneous polynomial of positive degree on
\(\R^N\), \(N\ge2\), satisfying \eqref{eq:harmonic-minimal-weight}, and set
\(U=\{S>0\}\).  Let \(h\) be smooth on \(U\), put \(a=e^h\),
and suppose that a smooth function \(\tau\) on \(U\) satisfies
\begin{equation}\label{eq:harmonic-log-compatibility}
 \Gamma(F,h)=\tau F.
\end{equation}
Thus \(\nabla h\) is tangent to the regular zero set. Introduce
\begin{equation}\label{eq:harmonic-log-auxiliary}
 H_h=\Gamma(h,h),\qquad Q_h=H_h+2\tau,\qquad
 P_h=Q_h\Delta h-\tfrac12\Gamma(h,Q_h),
\end{equation}
and the boundary coefficient
\begin{equation}\label{eq:harmonic-log-D0}
 D_h=S(\Delta h-\tau-H_h)-q-\tfrac12\Gamma(h,S).
\end{equation}
These functions are defined on \(U\) without division by \(F\).

\begin{lemma}[Harmonic multiplier identity]
\label{lem:harmonic-log-multiplier}
With the preceding notation, \(G=aF\) satisfies
\begin{align}
 |\nabla G|^2&=a^2(S+F^2Q_h),
 \label{eq:harmonic-log-gradient}\\
 \cL(G)&=a^3F\left(D_h+F^2P_h-
                    \tfrac F2\Gamma(F,Q_h)\right).
 \label{eq:harmonic-log-defect}
\end{align}
Along \(C_{\reg}=\{F=0,S>0\}\), one has
\begin{equation}\label{eq:harmonic-log-Jacobi}
 \frac{D_h}{S}=-\frac{J_Cw}{w},\qquad
 w=(a\sqrt S)^{-1}.
\end{equation}
If in addition \(\Gamma(F,Q_h)=\kappa F\) for a smooth function
\(\kappa\), then \eqref{eq:harmonic-log-defect} becomes
\begin{equation}\label{eq:harmonic-log-two-coefficients}
 \cL(G)=a^3F(D_h+F^2D_{1,h}),\qquad
 D_{1,h}=Q_h\Delta h-\tfrac12\kappa-
                         \tfrac12\Gamma(h,Q_h).
\end{equation}
\end{lemma}

\begin{proof}
The product and chain rules give
\[
 \nabla G=a(\nabla F+F\nabla h),\qquad
 \Delta G=aF(\Delta h+Q_h).
\]
The first formula gives \eqref{eq:harmonic-log-gradient}.  Set
\(V=S+F^2Q_h\).  Differentiating \(|\nabla G|^2=a^2V\) yields
\[
 a^{-3}\Gamma(G,|\nabla G|^2)
 =2FV(\tau+H_h)+\Gamma(F,V)+F\Gamma(h,V).
\]
Here
\[
 \begin{aligned}
 \Gamma(F,V)&=2qF+2FSQ_h+F^2\Gamma(F,Q_h),\\
 \Gamma(h,V)&=\Gamma(h,S)+2\tau F^2Q_h+
                                      F^2\Gamma(h,Q_h).
 \end{aligned}
\]
Substitute these identities in
\(\cL(G)=|\nabla G|^2\Delta G-
\tfrac12\Gamma(G,|\nabla G|^2)\), and use
\(Q_h=H_h+2\tau\).  The terms proportional to \(F\) and
\(F^3\) give \eqref{eq:harmonic-log-defect}.
The additional divisibility assumption gives
\eqref{eq:harmonic-log-two-coefficients}.

On the regular zero set, \(\cK_F(a)=a^3D_h\), so
\eqref{eq:boundary-potential-algebraic} gives
\eqref{eq:harmonic-log-Jacobi}.
\end{proof}

\begin{theorem}[A monotonicity criterion for harmonic multipliers]
\label{thm:harmonic-monotonicity-criterion}
In the setting of \cref{lem:harmonic-log-multiplier}, suppose on
\(U\) that
\begin{equation}\label{eq:harmonic-log-monotonicity}
 Q_h\ge0,\qquad D_h\ge0,\qquad \Delta h\ge0,\qquad
 F\Gamma(F,Q_h)\le0,\qquad \Gamma(h,Q_h)\le0.
\end{equation}
Then \(\nabla(aF)\ne0\) on \(U\), and
\(F\cL(aF)\ge0\) there. With
\(Z=\R^N\setminus U=\{\nabla F=0\}\), the field
\[
 X=\frac{\nabla(aF)}{|\nabla(aF)|}\quad\text{on }U
\]
is a two-sided subcalibration of \(E=\{F<0\}\) across \(Z\).
Consequently \(E\) is locally perimeter minimizing and
\(\partial\llbracket E\rrbracket\) is locally mass minimizing.
No regularity of the spherical link is assumed.
\end{theorem}

\begin{proof}
By \eqref{eq:harmonic-log-gradient},
\(|\nabla(aF)|^2\ge a^2S>0\).  In the bracket in
\eqref{eq:harmonic-log-defect}, each of
\(D_h\), \(F^2Q_h\Delta h\),
\(-F^2\Gamma(h,Q_h)/2\), and
\(-F\Gamma(F,Q_h)/2\) is nonnegative.  Thus
\(F\cL(aF)\ge0\), and
\eqref{eq:div-normalized-gradient} gives the required divergence
signs.  At every regular zero, \(\nabla(aF)=a\nabla F\), so
\(X\) is the exterior unit normal to \(E\).

By \cref{rem:polynomial-finite-perimeter,lem:square-free-critical-set},
\(E\) has locally finite perimeter and its algebraic critical set
\(Z\subset\{F=0\}\) has locally zero \(\cH^{N-1}\)-measure.
Apply \cref{thm:removable-subcalibration} to \(X\); no extension
of \(a\) across \(Z\) is needed.
\end{proof}

\begin{remark}[Alternative coefficient tests]
\label{rem:harmonic-alternative-tests}
If \(\Gamma(F,Q_h)=\kappa F\), the monotonicity conditions may be replaced by
\(Q_h\ge0\), \(D_h\ge0\), and \(D_{1,h}\ge0\).
Without this divisibility assumption, the conditions
\[
 Q_h\ge0,\quad D_h\ge0,\quad P_h\ge0,\quad
 \Gamma(F,Q_h)^2\le16D_hP_h
\]
also suffice: the last inequality gives
\(\tfrac12|F\Gamma(F,Q_h)|\le D_h+F^2P_h\).
All inequalities must hold throughout \(U\).
\end{remark}

For the natural choice \(a=\sqrt S\),
the quantities in the criterion use only \(q,S\) and their derivatives. Writing
\(K=\Delta S\) and \(T=\Gamma(S,S)\), one obtains
\begin{equation}\label{eq:harmonic-gradient-multiplier-data}
 \begin{aligned}
 h&=\tfrac12\log S,& \tau&=q/S,&
 Q_h&=\frac{T}{4S^2}+\frac{2q}{S},\\
 \Delta h&=\frac{K}{2S}-\frac{T}{2S^2},&&&
 D_h&=\frac K2-2q-\frac TS.
 \end{aligned}
\end{equation}
Under the spectral identities below, the two monotonicity signs hold,
and the remaining condition reduces to a lower bound for \(\Delta S\).

\begin{remark}[The nonharmonic correction]
\label{rem:nonharmonic-log-multiplier}
Suppose instead that \(\cL(F)=-qF\), without assuming
\(\Delta F=0\), and retain \(U=\{S>0\}\),
\eqref{eq:harmonic-log-compatibility}, and the auxiliary functions above.
The term \(2S\Delta F\) in \(\Gamma(F,S)\) changes only the defect:
\[
 \cL(aF)=a^3F\left(D_h+FQ_h\Delta F+F^2P_h
                         -\tfrac F2\Gamma(F,Q_h)\right).
\]
The gradient identity is unchanged, so the square-completion test of
\cref{rem:harmonic-alternative-tests} becomes
\[
 Q_h\ge0,\qquad D_h\ge0,\qquad P_h\ge0,\qquad
 \bigl(\Gamma(F,Q_h)-2Q_h\Delta F\bigr)^2\le16D_hP_h.
\]
It gives comparison when \(Z=\{S=0\}\) is locally
\(\cH^{N-1}\)-null, in particular for square-free homogeneous \(F\)
by \cref{lem:square-free-critical-set}. The natural choice
\(h=\frac12\log S\) must still satisfy the compatibility assumption:
now \(\Gamma(F,h)=\Delta F+qF/S\), which need not be divisible
smoothly by \(F\) across its regular zero set.
\end{remark}


\subsection{Radial multipliers in arbitrary degree}
\label{sec:general-radial-multipliers}

Let \(F\) be a nonzero homogeneous polynomial of degree \(e\ge2\),
put \(S=|\nabla F|^2\), and suppose \(\cL(F)=-qF\) for a
polynomial \(q\).
The identity in \cref{prop:radial-master-identity} gives the following
radial specialization.
For \(\sigma>0\), set
\begin{align}
 A_\sigma&=\sigma(N-1-2e-\sigma),\qquad
 c_\sigma=\sigma(\sigma+2e),\notag\\
 B_\sigma&=c_\sigma\bigl(e+\sigma(N-1)\bigr).
 \label{eq:general-radial-coefficients}
\end{align}
Then, on \(r>0\),
\begin{align}
 |\nabla(r^\sigma F)|^2
 &=r^{2\sigma}\left(S+c_\sigma\frac{F^2}{R}\right),
 \label{eq:general-radial-gradient}\\
 \cL(r^\sigma F)&=r^{3\sigma-2}F\,\mathscr H_\sigma,
 \label{eq:general-radial-defect}\\
 \mathscr H_\sigma
 &=A_\sigma S-Rq+c_\sigma F\Delta F
                  +B_\sigma\frac{F^2}{R}.
 \label{eq:general-radial-bracket}
\end{align}
The exact sign condition is \(\mathscr H_\sigma\ge0\) on \(r>0\).
A sufficient test follows.

\begin{proposition}[A radial criterion without harmonicity]
\label{prop:general-radial-existence}
Let \(F\) be a homogeneous polynomial of degree \(e\ge2\) on
\(\R^N\), \(N\ge3\), taking both signs,
and suppose \(\cL(F)=-qF\) for a polynomial \(q\).
Assume that \(Z=\{\nabla F=0\}\) has dimension at most \(N-2\),
as holds whenever \(F\) is square-free by
\cref{lem:square-free-critical-set}.
If constants \(\kappa,\eta\ge0\) and \(\sigma>0\) satisfy
\begin{equation}\label{eq:general-radial-comparisons}
 Rq\le\kappa S,\qquad F\Delta F\ge-\eta S
 \quad\hbox{on }\R^N,
\end{equation}
and
\begin{equation}\label{eq:general-radial-sufficient}
 \sigma\bigl[N-1-2e(1+\eta)-(1+\eta)\sigma\bigr]\ge\kappa,
\end{equation}
then \(G=r^\sigma F\) generates a two-sided subcalibration and
\(\partial\llbracket\{F<0\}\rrbracket\) is strictly area minimizing.
Moreover, on \(C_{\reg}=\{F=0,S>0\}\),
\begin{equation}\label{eq:general-radial-Hardy}
 Q_C(u)\ge
 \left[\frac{(N-1-2e)^2}{4}-\kappa\right]
 \int_Cr^{-2}u^2,
 \qquad u\in C_c^\infty(C_{\reg}).
\end{equation}
\end{proposition}

\begin{proof}
The assumptions give
\[
 \mathscr H_\sigma\ge
 (A_\sigma-c_\sigma\eta-\kappa)S
 +B_\sigma\frac{F^2}{R}\ge0.
\]
Since \(\sigma,c_\sigma,B_\sigma>0\), the bracket is positive off
the cone. Euler's identity and \eqref{eq:general-radial-gradient}
give \(Z\subset\{F=0\}\) and \(\nabla G\ne0\) off \(Z\).
The continuous potential has the signs of \(F\), and the algebraic
set \(Z\) is locally \(\cH^{N-1}\)-null. Thus
\cref{cor:gradient-subcalibration,thm:strict-minimality-flux}
give minimization and its strict form.

The Jacobi trace \eqref{eq:Jacobi-trace} and
\eqref{eq:general-radial-defect} give, on \(C_{\reg}\),
\begin{equation}\label{eq:general-radial-Jacobi}
 -\frac{J_Cw}{w}=\frac{A_\sigma}{R}-\frac qS
 \ge\frac{A_\sigma-\kappa}{R},
 \qquad w=r^{-\sigma}S^{-1/2}.
\end{equation}
Apply \cref{thm:ground-state-gap} with \(\ell=e+\sigma\) and use
\[
 \left(e+\sigma-\frac{N-1}{2}\right)^2+A_\sigma
 =\frac{(N-1-2e)^2}{4}
\]
to obtain \eqref{eq:general-radial-Hardy}.
\end{proof}

Writing \(b_\eta=N-1-2e(1+\eta)\), the left side of
\eqref{eq:general-radial-sufficient} is maximized, when
\(b_\eta>0\), at \(\sigma=b_\eta/[2(1+\eta)]\).
This yields the explicit sufficient dimension condition
\begin{equation}\label{eq:general-radial-dimension}
 b_\eta>0,\qquad
 N\ge1+2e(1+\eta)+2\sqrt{(1+\eta)\kappa}.
\end{equation}
Separating the terms in \eqref{eq:general-radial-bracket} can lose sharpness:

\begin{example}[Nonsharpness of the simplified criterion]
\label{ex:radial-separated-nonsharp}
Let \(F=7u-v\), where \(u=|x|^2\), \(v=|y|^2\), and
\((x,y)\in\R^2\times\R^8\). This is a classical Lawson cone;
see \cite{Liu19} for quantitative subcalibrations in this family.
Here
\[
 \Delta F=12,\qquad S=196u+4v,\qquad \cL(F)=-56F.
\]
The least nonnegative constants in
\eqref{eq:general-radial-comparisons} are
\(\kappa=14\) and \(\eta=3\), both attained at \(u=0<v\).
Consequently \eqref{eq:general-radial-sufficient} would require
\(\sigma(-7-4\sigma)\ge14\), which fails for every
\(\sigma>0\). Nevertheless, the radial polynomial potential
\(G=RF\) satisfies the exact identity
\begin{equation}\label{eq:radial-Lawson-counterexample}
 \cL(G)=64RF\bigl[(v-11u)^2+96u^2\bigr].
\end{equation}
Since \(S>0\) away from the origin, this gives a two-sided
subcalibration with strict signed divergence off the cone.
Thus \eqref{eq:general-radial-sufficient} need not hold for any
positive exponent, even with the least admissible comparison constants.
\end{example}

Finite comparison constants also need not exclude nonzero critical
points without harmonicity. For this explicitly reducible comparison,
outside the introductory convention, put \(u=|x|^2\), \(v=|y|^2\), with
\((x,y,z)\in\R^4\times\R^4\times\R^2\), the square-free
but reducible polynomial \(F=u^2-v^2\) satisfies
\[
 S=16(u^3+v^3),\qquad \Delta F=24(u-v),\qquad
 \cL(F)=192F(u-v)^2.
\]
Hence \(\kappa=\eta=0\) are admissible in
\eqref{eq:general-radial-comparisons}, and \(\sigma=1/2\)
satisfies \eqref{eq:general-radial-sufficient}. Its zero set is the
Simons cone in \(\R^8\) times \(\R^2\), with critical set
\(\{x=y=0\}\). The quartic defining equation includes the
positive factor \(u+v\) away from that set.

If \(F\) belongs to a Laplacian algebra \(\cA\) containing
\(R\), then \(S,\Delta F\in\cA\), and the division property
of \cref{prop:MR-finite-and-division} gives \(q\in\cA\).
Thus \(R\mathscr H_\sigma\in\cA\), and its sign on \(R>0\)
can be tested on the real quotient using the closure table.

\begin{corollary}[A harmonic radial criterion in arbitrary degree]
\label{cor:harmonic-radial-existence}\label{main:harmonic-radial}
Under the hypotheses of \cref{prop:harmonic-power-Jacobi}, suppose that
for a constant \(\kappa\ge0\),
\begin{equation}\label{eq:harmonic-radial-bound}
 Rq\le\kappa S\quad\hbox{on }\R^N,\qquad
 N\ge2e+1+2\sqrt\kappa.
\end{equation}
Then, with \(\sigma=(N-1-2e)/2\), the potential \(G=r^\sigma F\)
generates a two-sided subcalibration. The phase \(\{F<0\}\) is locally
perimeter minimizing, and its oriented boundary current is locally
mass minimizing with support \(C=\{F=0\}\). The cone is strictly
area minimizing, and
\begin{equation}\label{eq:harmonic-radial-Hardy}
 Q_C(u)\ge\left[\frac{(N-1-2e)^2}{4}-\kappa\right]
 \int_Cr^{-2}u^2,
 \qquad u\in C_c^\infty(C\setminus\{0\}).
\end{equation}
\end{corollary}

\begin{proof}
Write \(b=N-1-2e\). By
\cref{prop:harmonic-global-comparison} below,
\(S>0\) off the origin and \(\kappa\ge e(e-1)>0\).
The dimension assumption gives \(b>0\), and
\(\sigma=b/2\) satisfies
\(\sigma(b-\sigma)=b^2/4\ge\kappa\).
Harmonicity gives both nonempty phases and
\(F\Delta F=0\). Thus \cref{prop:general-radial-existence}
applies with \(\eta=0\) and \(Z=\{0\}\), proving the
subcalibration, strict minimality, and Hardy bound.
The support of the boundary current is all of \(\{F=0\}\)
by \cref{rem:polynomial-finite-perimeter}.
\end{proof}

Any \(\sigma\ge0\) with \(\sigma(N-1-2e-\sigma)\ge\kappa\)
works; even integers give polynomial potentials, as in the
\hyperref[par:Simons-degree-example]{Simons example}. A local radial
subcalibration already implies strict stability for regular harmonic cones:

\begin{corollary}[Strict stability from a harmonic radial potential]
\label{cor:harmonic-radial-strict-stability}
Let $F$ be a nonzero real homogeneous harmonic polynomial
of degree $e\ge2$ on $\R^{n+1}$, $n\ge3$. Suppose $C=\{F=0\}$
is minimal, $\nabla F\ne0$ on $C\setminus\{0\}$, and its link
is nonempty, smooth, compact, and connected. If for some
$\alpha\in\R$ the potential $G=r^\alpha F$ satisfies
\[
 \operatorname{sgn}(F)\diver X_G\ge0\qquad\text{where }F\ne0
\]
in a two-sided neighborhood of $C\setminus\{0\}$ on which $X_G$
is defined, then $C$ is strictly stable.
\end{corollary}

\begin{proof}
The hypotheses already imply that $F$ is real irreducible.
Indeed, \cref{lem:square-free-critical-set} shows that $F$ is
square-free and each real irreducible factor has a nonempty zero
set on the unit sphere. These closed sets are pairwise disjoint,
since a common zero of two factors would give $\nabla F=0$.
Their union is the connected link, so there can be only one factor.
Set $R=r^2$ and $S=|\nabla F|^2$.
By \cref{prop:divisibility}, $\cL(F)=-qF$ for a homogeneous
polynomial $q$ of degree $2e-4$, possibly zero.
The sign hypothesis and \cref{thm:Jacobi-trace} make
$w=r^{-\alpha}S^{-1/2}$ a positive Jacobi supersolution.
Suppose that the Hardy gap is zero. Equation
\eqref{eq:critical-boundary-rigidity} gives
$e+\alpha=n/2$ and $J_Cw=0$.
With $c=\alpha(\alpha+2e)(e+\alpha n)$, the radial identity
of \cref{prop:radial-master-identity} and $\Delta F=0$ give
\begin{equation}\label{eq:harmonic-critical-radial-defect}
 \cL(G)=r^{3\alpha-4}FB,\qquad
 B=R(\alpha^2S-Rq)+cF^2.
\end{equation}
Thus $B\ge0$ near the cone, and $J_Cw=0$ together with
\cref{thm:Jacobi-trace} gives $B|_{C\setminus\{0\}}=0$.
The divisibility argument in \cref{prop:divisibility} yields
$B=FB_1$. At a regular zero, $F$ changes sign, so $B\ge0$
on both sides forces $B_1$ to vanish there as well.
Applying the same argument again gives $F^2\mid B$.
Since $F$ has nonzero real zeros, it is coprime to $R$.
Equation \eqref{eq:harmonic-critical-radial-defect} therefore implies
$F^2\mid\alpha^2S-Rq$. Its degree is at most $2e-2<2e$, so
\begin{equation}\label{eq:harmonic-critical-radial-obstruction}
 Rq=\alpha^2S\qquad\text{identically}.
\end{equation}
Choose an extremum $p$ of $F$ on the unit sphere with $F(p)\ne0$.
Euler's identities give
$\nabla F(p)=eF(p)p$ and
$\Hess F(p)p=(e-1)\nabla F(p)$.
Using harmonicity and $\cL(F)=-qF$ gives
$q(p)=e(e-1)S(p)$, with $S(p)>0$.
Consequently \eqref{eq:harmonic-critical-radial-obstruction}
forces $\alpha^2=e(e-1)$.
Together with $\alpha=n/2-e$, this becomes $n^2=4e(n-1)$,
which is impossible for integers $n\ge3$, since it would imply
$n-1\mid1$.
\end{proof}

\begin{proposition}[Global comparison and the critical set]
\label{prop:harmonic-global-comparison}
Let \(F\not\equiv0\) be a homogeneous polynomial of degree \(e\ge2\)
on \(\R^N\), \(N\ge3\), satisfying
\(\Gamma(F,S)=2qF\), where \(S=|\nabla F|^2\) and \(q\) is a
polynomial. Then
\begin{equation}\label{eq:global-comparison-regularity}
 \exists\,\kappa<\infty:\ Rq\le\kappa S
 \quad\Longleftrightarrow\quad
 \{\nabla F=0\}=\{0\}.
\end{equation}
In this case the least comparison constant is
\begin{equation}\label{eq:harmonic-optimal-comparison}
 \kappa_F=\max_{\Sph^{N-1}}\frac qS\ge e(e-1).
\end{equation}
Consequently the hypotheses of
\cref{cor:harmonic-radial-existence} require \(N\ge4e\).
\end{proposition}

\begin{proof}
If \(S>0\) off the origin, compactness of the unit sphere and
homogeneity give the finite maximum in
\eqref{eq:harmonic-optimal-comparison}.
Conversely, suppose \(c\ne0\) and \(\nabla F(c)=0\).
Euler's identity gives \(F(c)=0\). Write
\[
 F(c+y)=P_m(y)+O(|y|^{m+1}),\qquad m\ge2,
\]
where \(P_m\not\equiv0\) is the first nonzero homogeneous Taylor
term. Choose a unit vector \(v\) maximizing \(|P_m|\), and put
\(p=P_m(v)\ne0\). Then
\(\nabla P_m(v)=mpv\) and
\(\Hess P_m(v)(v,v)=m(m-1)p\). Along \(c+tv\),
\begin{align*}
 S&=m^2p^2t^{2m-2}(1+O(t)),\\
 q=\frac{\Hess F(\nabla F,\nabla F)}{F}
  &=m^3(m-1)p^2t^{2m-4}(1+O(t)).
\end{align*}
It follows that
\begin{equation}\label{eq:critical-comparison-blowup}
 \frac{Rq}{S}(c+tv)
 =\frac{|c|^2m(m-1)}{t^2}+O(t^{-1})\longrightarrow+\infty.
\end{equation}
This rules out a finite comparison constant, including when
\(q(c)=0\).

For the lower bound, take a unit vector maximizing \(|F|\).
The same calculation, now with degree \(e\) at the origin, gives
\(q/S=e(e-1)\) at that vector. Finally, \(4e-1<2e+1+2\sqrt{e(e-1)}<4e\),
so the integer dimension in the harmonic radial criterion is at least \(4e\).
\end{proof}

\begin{proposition}[Rigidity at the sharp comparison constant]
\label{prop:harmonic-sharp-comparison}
Under the algebraic hypotheses of
\cref{prop:harmonic-global-comparison}, let
\(M=\max_{\Sph^{N-1}}|F|\). Then
\begin{equation}\label{eq:harmonic-comparison-rigidity}
 Rq\le e(e-1)S
 \quad\Longleftrightarrow\quad
 S=e^2M^2R^{e-1}.
\end{equation}
If also \(\Delta F=0\) and \(N\ge3\), the equality case is
precisely a harmonic Cartan--M\"unzner polynomial after normalization
by \(M\).
\end{proposition}

\begin{proof}
Put \(f=F|_{\Sph^{N-1}}\) and \(s=S|_{\Sph^{N-1}}\).
The restriction of \(F\) to any great circle is a real trigonometric
polynomial of degree at most \(e\). Applying the classical
Bernstein--Szeg\H{o} inequality
\cite[Section~5.1, Theorem~5.1.3]{BorweinErdelyi95} to a great circle in the
direction of \(\nabla_{\Sph}f\) gives
\begin{equation}\label{eq:spherical-Bernstein-bound}
 s=|\nabla_{\Sph}f|^2+e^2f^2\le e^2M^2.
\end{equation}
On the other hand, homogeneity and \(\Gamma(F,S)=2qF\) give
\[
 \Gamma_{\Sph}(f^2,s)
 =4f^2\bigl(q-e(e-1)s\bigr)\le0.
\]
Thus \(s\) is nonincreasing along the gradient ascent flow of
\(f^2\). The critical set of \(f^2\) is semialgebraic with
finitely many connected components, and \(f^2\) is constant on each
component. Choose \(a<M^2\) above every other critical value.
Every forward trajectory starting in the nonempty open set
\(\{f^2>a\}\) has its accumulation points in the critical set at
value \(M^2\): the value increases along the trajectory, and every
accumulation point of a gradient flow on this compact sphere is
critical. At those points \(s=e^2M^2\). Monotonicity of \(s\)
and \eqref{eq:spherical-Bernstein-bound} therefore force
\(s=e^2M^2\) throughout \(\{f^2>a\}\). Homogeneity and the
polynomial identity principle give the asserted formula for \(S\)
on all of \(\R^N\).

Conversely, differentiating that formula yields
\(q=e^3(e-1)M^2R^{e-2}\), proving the comparison with equality.
When \(F\) is harmonic, \(F/M\) satisfies both Cartan--M\"unzner
equations with equal multiplicities.
\end{proof}

\begin{lemma}[Degree of a nonnegative polynomial first integral]
\label{lem:nonnegative-first-integral-degree}
Let \(F\not\equiv0\) be a homogeneous polynomial of degree \(e\ge1\)
on \(\R^N\), \(N\ge2\), with
\(\nabla F\ne0\) on \(\{F=0\}\setminus\{0\}\).
If \(K\not\equiv0\) is a homogeneous polynomial of degree \(\ell>0\),
\(K\ge0\), and \(\Gamma(F,K)=0\), then \(\ell\ge2e\).
\end{lemma}

\begin{proof}
Choose \(p\in\Sph^{N-1}\) with
\(K(p)=B=\max_{\Sph^{N-1}}K>0\).
The constrained maximum and Euler's identity give
\(\nabla K(p)=\ell Bp\), hence
\[
 0=\Gamma(F,K)(p)=e\ell BF(p).
\]
Thus \(F(p)=0\) and \(w=\nabla F(p)\ne0\), with \(w\perp p\).
Differentiating Euler's identity gives
\(\Hess F(p)p=(e-1)w\). Differentiating
\(\Gamma(F,K)=0\) at \(p\) then gives
\[
 \Hess K(p)w=-\Hess F(p)\nabla K(p)=-\ell B(e-1)w.
\]
For \(v=w/|w|\), the function
\(T(t)=K(p\cos t+v\sin t)\) is a real trigonometric polynomial
of degree at most \(\ell\), with \(0\le T\le B\), and
\[
 T''(0)=\Hess K(p)(v,v)-\nabla K(p)\cdot p=-e\ell B.
\]
The trigonometric Bernstein inequality, applied twice to \(T-B/2\)
\cite[Section~5.1, Theorem~5.1.4]{BorweinErdelyi95}, yields
\[
 e\ell B\le\|T''\|_\infty
 \le\ell^2\|T-B/2\|_\infty\le\frac{\ell^2B}{2}.
\]
Consequently \(\ell\ge2e\).
\end{proof}

\begin{proposition}[Rigidity for a radial minimality weight]
\label{prop:regular-radial-weight-rigidity}
Let \(F\not\equiv0\) be a homogeneous polynomial of degree \(e\ge2\) on
\(\R^N\), \(N\ge3\), with \(\operatorname{Crit}(F)=\{0\}\).
Put \(S=|\nabla F|^2\) and
\(M_F=\max_{\Sph^{N-1}}|F|\). If
\[
 \Gamma(F,S)=2\lambda R^{e-2}F
\]
for a constant \(\lambda\), then
\[
 \lambda=e^3(e-1)M_F^2,\qquad
 S=e^2M_F^2R^{e-1}=\frac{\lambda}{e(e-1)}R^{e-1}.
\]
If also \(\Delta F=0\), then \(F/M_F\) is a harmonic
Cartan--M\"unzner polynomial. In particular,
\[
 \cA_R(F)=\R[R,F],\qquad e\in\{2,3,4,6\},\qquad
 \kappa_F=e(e-1).
\]
\end{proposition}

\begin{proof}
After replacing \(F\) by \(-F\) if necessary, choose
\(p\in\Sph^{N-1}\) with \(F(p)=M_F\). Then
\(\nabla F(p)=eM_Fp\) and \(S(p)=e^2M_F^2\).
Since \(S\) has degree \(2e-2\),
\[
 2\lambda M_F=\Gamma(F,S)(p)
 =eM_F(2e-2)S(p)=2e^3(e-1)M_F^3.
\]
This determines \(\lambda\). The spherical Bernstein--Szeg\H{o}
bound \eqref{eq:spherical-Bernstein-bound}, extended by homogeneity,
shows that \(K:=e^2M_F^2R^{e-1}-S\ge0\).
Using \(\Gamma(F,R)=2eF\), we obtain
\[
 \Gamma(F,K)
 =2\bigl[e^3(e-1)M_F^2-\lambda\bigr]R^{e-2}F=0.
\]
If \(K\not\equiv0\), its degree \(2e-2\) contradicts
\cref{lem:nonnegative-first-integral-degree}. Hence \(K=0\).

If \(\Delta F=0\), the normalized equations
\(|\nabla(F/M_F)|^2=e^2R^{e-1}\) and \(\Delta(F/M_F)=0\)
are the harmonic Cartan--M\"unzner equations, as in
\cref{prop:harmonic-sharp-comparison}. Together with
\(\Gamma(R,F)=2eF\) and \(\Delta R=2N\), they show that
\(\cA_R(F)=\R[R,F]\). The nonempty zero link is regular by hypothesis
and minimal since \(\cL(F)=-\lambda R^{e-2}F\); hence
\cref{thm:rank-one-classification} gives the degree restriction.
Finally, \(Rq/S=e(e-1)\) off the origin.
\end{proof}

\begin{example}[Harmonic isoparametric polynomials]
\label{ex:harmonic-isoparametric-radial}
For a normalized harmonic Cartan--M\"unzner polynomial of degree
\(e\ge2\) on \(\R^N\), \(N\ge3\),
\[
 S=e^2R^{e-1},\qquad q=e^3(e-1)R^{e-2},\qquad
 \kappa_F=e(e-1).
\]
The dimension condition of \cref{cor:harmonic-radial-existence}
is therefore equivalent to \(N\ge4e\). In this range the choice
\(\sigma=2\lfloor e/2\rfloor\in\{e-1,e\}\) satisfies
\(\sigma(N-1-2e-\sigma)\ge e(e-1)\), so
\[
 G=R^{\lfloor e/2\rfloor}F
\]
is a polynomial subcalibration potential.
The realizable equal-multiplicity parameters satisfying this threshold
are listed in \cref{sec:harmonic-isoparametric-scope}, together with
the Clifford-algebra explanation for the quartic restriction.
By \cref{prop:regular-radial-weight-rigidity}, they exhaust
\cref{cor:harmonic-radial-existence} when \(q\) is radial:
harmonic quadrics in even dimensions \(N\ge8\), and Cartan cubics
in dimensions \(14\) and \(26\).
\end{example}

\begin{proposition}[Positive multipliers at a harmonic critical zero]
\label{prop:harmonic-positive-multiplier-obstruction}
\label{rem:harmonic-radial-singular-obstruction}
Let \(F\not\equiv0\) be a harmonic polynomial, and suppose
\(F(c)=0=\nabla F(c)\). If \(a\) is \(C^2\) near \(c\) and
\(a(c)>0\), then \(G=aF\) cannot satisfy \(G\cL(G)\ge0\)
throughout any neighborhood of \(c\).
\end{proposition}

\begin{proof}
Use the first nonzero Taylor term \(P_m\), the vector \(v\), and
\(p=P_m(v)\ne0\) as in the proof of
\cref{prop:harmonic-global-comparison}. Harmonicity gives
\(\Delta(aF)=2\Gamma(a,F)+F\Delta a=O(t^{m-1})\) along
\(c+tv\). The leading Hessian term therefore yields
\[
 F(c+tv)\cL(aF)(c+tv)
 =-a(c)^3m^3(m-1)p^4t^{4m-4}+o(t^{4m-4})<0
\]
for all sufficiently small positive \(t\). Here \(a>0\) and
\(\nabla(aF)\ne0\), so the required sign fails.
\end{proof}

\begin{remark}[Quartic and quintic scope]
\label{rem:higher-degree-radial-scope}
The criterion of \cref{cor:harmonic-radial-existence} applies in every
degree but requires a regular link. In degrees four and five its
conditions read, respectively, \(N\ge16\),
\(\kappa_F\le(N-9)^2/4\), and \(N\ge20\),
\(\kappa_F\le(N-11)^2/4\). Degree five is absent from the
Cartan--M\"unzner list, so \cref{prop:harmonic-sharp-comparison}
forces \(\kappa_F>20\) in that case.
Moreover, \cref{prop:regular-radial-weight-rigidity} forces \(q\)
to be nonradial in any degree-five example satisfying the criterion.
The determinant cones in \cref{main:complex-determinants} instead
have nonzero critical points and use nonradial minor-norm multipliers.
Likewise, adding cylindrical directions cannot produce new examples
of the radial criterion, by
\cref{prop:harmonic-positive-multiplier-obstruction}.
\end{remark}

\subsection{Homogeneous potentials on a spherical quotient}
\label{sec:spherical-quotient-potentials}

Let \(I=(I_1,\ldots,I_j)\), where each invariant \(I_i\) is
homogeneous of degree \(d_i\), and write
\[
 z_i=\frac{I_i}{r^{d_i}},\qquad
 K_{\Sph}=z(\Sph^{N-1}).
\]
For an ambient closure containing \(R\), Euler's identity
\(\nabla_{\Sph}I_i=\nabla I_i-d_iI_i\theta\) and the polar
Laplacian give the spherical closure data
\begin{equation}\label{eq:spherical-closure-from-ambient}
 g^{\Sph}_{ij}=\left.\Gamma(I_i,I_j)\right|_{R=1}-d_id_jz_iz_j,
 \qquad
 b^{\Sph}_i=\left.\Delta I_i\right|_{R=1}
              -d_i(d_i+N-2)z_i.
\end{equation}
More generally, assume the spherical data close as functions of \(z\).
Let \(\widehat\Gamma_{\Sph}\), \(\widehat\Delta_{\Sph}\) be their
operators from \cref{thm:finite-reduction}, with signs tested on
the actual image \(K_{\Sph}\).

\begin{proposition}[The spherical equation for a homogeneous potential]
\label{prop:spherical-quotient-potential}
Let $\varphi$ be smooth on a neighborhood of the relevant part of
$K_{\Sph}$ in the ambient invariant-coordinate space.
For \(G=r^\ell \varphi(z)\), define
\begin{align}
 W_\ell[\varphi]&=\ell^2\varphi^2+\widehat\Gamma_{\Sph}(\varphi,\varphi),
 \label{eq:general-spherical-W}\\
 \mathfrak N_\ell[\varphi]
 &=W_\ell[\varphi]\bigl(\widehat\Delta_{\Sph}\varphi+(N-1)\ell \varphi\bigr)
       -\tfrac12\widehat\Gamma_{\Sph}(\varphi,W_\ell[\varphi]).
 \label{eq:general-spherical-N}
\end{align}
Then
\begin{equation}\label{eq:general-spherical-L}
 |\nabla G|^2=r^{2\ell-2}W_\ell[\varphi],\qquad
 \cL(G)=r^{3\ell-4}\mathfrak N_\ell[\varphi].
\end{equation}
Where \(W_\ell[\varphi]>0\), one has
\(r\diver X_G=\mathfrak N_\ell[\varphi]/W_\ell[\varphi]^{3/2}\).
These identities hold at singular quotient values whenever their
pullbacks are smooth and \(W_\ell[\varphi]>0\).
\end{proposition}

\begin{proof}
Polar differentiation gives
\[
 \nabla G=r^{\ell-1}(\ell \varphi\,\partial_r+\nabla_{\Sph}\varphi),\qquad
 \Delta G=r^{\ell-2}
       [\widehat\Delta_{\Sph}\varphi+\ell(\ell+N-2)\varphi].
\]
The radial and angular terms give
\[
 \tfrac12\Gamma(G,|\nabla G|^2)
 =r^{3\ell-4}\left[
  \ell(\ell-1)\varphi W_\ell[\varphi]
  +\tfrac12\widehat\Gamma_{\Sph}(\varphi,W_\ell[\varphi])\right].
\]
Subtracting this expression from \(|\nabla G|^2\Delta G\) proves the formula.
\end{proof}

For a nonzero homogeneous polynomial \(F\) of degree \(e\geq1\),
put \(R=r^2\) and \(S=|\nabla F|^2\).
When these quantities give a closed quotient as above, the angular
variables are
\begin{equation}\label{eq:rank-two-general-variables}
 f=\frac{F}{R^{e/2}},\qquad
 s=\frac{S}{R^{e-1}},\qquad K_{\Sph}=\{(f(x),s(x)):r=1\}.
\end{equation}
If this image has dimension two, a homogeneous invariant potential
with a positive multiplier has the form
\[
 G=r^\ell f\,a(f,s),\qquad a>0.
\]
The power family \eqref{eq:power-family} corresponds to
\(\varphi(f,s)=fs^\gamma\) and
\(\ell=e+2\alpha+2\gamma(e-1)\).  It is a restricted family of
angular multipliers.

\begin{remark}[Phase-dependent quotient potentials]
\label{cor:rank-two-paired-potentials}
Let \(N\geq3\).  Suppose \(E=\{F<0\}\) has smooth boundary off a
closed conical set \(Z\) containing the origin and satisfying the
null-set hypothesis of \cref{thm:paired-subcalibration}. Use the quotient
\eqref{eq:rank-two-general-variables}.
Choose degrees \(\ell_-\), \(\ell_+\) and put
\(h_\pm=f a_\pm(f,s)\), where the positive multipliers \(a_\pm\)
are smooth on invariant-coordinate neighborhoods of the respective
closed-phase images outside \(Z\).
Assume there that \(W_{\ell_\pm}[h_\pm]>0\), and assume
\[
 \mathfrak N_{\ell_-}[h_-]\leq0\quad(f<0),\qquad
 \mathfrak N_{\ell_+}[h_+]\geq0\quad(f>0).
\]
Then \(G_\pm=r^{\ell_\pm}h_\pm\) give
\eqref{eq:paired-weighted-comparison} with weights on the respective phases
\[
 q_\pm=\frac1r\,
  \frac{|\mathfrak N_{\ell_\pm}[h_\pm]|}
       {W_{\ell_\pm}[h_\pm]^{3/2}},
\]
extended by zero outside their phases. The boundary is area minimizing,
and strictly area minimizing if both unit-ball phase integrals are positive.
Indeed, on the regular boundary,
\(\nabla G_\pm=r^{\ell_\pm-e}a_\pm\nabla F\), so both normalized
fields equal \(\nu_E\).  The spherical equation gives their phase
signs and weights.  Apply \cref{rem:phase-local-fields},
\cref{thm:paired-subcalibration}, and
\cref{cor:paired-strict-minimality}. The two multipliers need not agree
on \(f=0\).
\end{remark}

A useful initial family is \(G=R^\alpha S^\gamma F\), with
\(S=|\nabla F|^2\).  Its full multiplier identities are collected in
\cref{app:power-identities}, especially
\cref{prop:combined-power-multiplier}.
